%====================================================
% PART A
% ISO/IEC 27001
%====================================================

\chapter{ISO/IEC 27001: Information Security Management}

This chapter examines ISO/IEC 27001:2022 as an internationally recognized framework for establishing, implementing, maintaining, and continually improving an Information Security Management System (ISMS). It explores the standard's structure, implementation methodology, and risk management practices through practical organizational scenarios.

%====================================================
\section{Foundational Concepts}
%====================================================

\subsection{Overview of ISO/IEC 27001}

%---------------------------------------

\subsection{Information Security Management System (ISMS)}

%---------------------------------------

\subsection{The CIA Triad}

\subsubsection{Confidentiality}

\subsubsection{Integrity}

\subsubsection{Availability}

%---------------------------------------

\subsection{ISO/IEC 27001 versus ISO/IEC 27002}

%---------------------------------------

\subsection{Business Benefits of ISO/IEC 27001 Certification}

%%%%%%%%%%%%%%%%%%%%%%%%%%%%%%%%%%%%%%%%%%%%%%%%%%%%%

\section{Structure and Security Controls}

%====================================================

\subsection{Annex A Control Themes}

\subsubsection{Organizational Controls}

\subsubsection{People Controls}

\subsubsection{Physical Controls}

\subsubsection{Technological Controls}

%---------------------------------------

\subsection{Selected Annex A Controls}

\subsubsection{Control One}

% Objective

% Risk Mitigated

% Practical Implementation

%---------------------------------------

\subsubsection{Control Two}

%---------------------------------------

\subsubsection{Control Three}

%---------------------------------------

\subsubsection{Control Four}

%---------------------------------------

\subsubsection{Control Five}

%---------------------------------------

\subsection{Statement of Applicability (SoA)}

%%%%%%%%%%%%%%%%%%%%%%%%%%%%%%%%%%%%%%%%%%%%%%%%%%%%%

\section{PDCA Cycle and ISMS Implementation}

%====================================================

\subsection{Plan}

%---------------------------------------

\subsection{Do}

%---------------------------------------

\subsection{Check}

%---------------------------------------

\subsection{Act}

%---------------------------------------

\subsection{Implementation Challenges}

\subsubsection{Challenge One}

\subsubsection{Challenge Two}

\subsubsection{Challenge Three}

%---------------------------------------

\subsection{Internal Audit Process}

%---------------------------------------

\subsection{Management Review and Continual Improvement}

%%%%%%%%%%%%%%%%%%%%%%%%%%%%%%%%%%%%%%%%%%%%%%%%%%%%%

\section{Risk Management}

%====================================================

\subsection{Risk Assessment Methodology}

\subsubsection{Risk Identification}

\subsubsection{Risk Analysis}

\subsubsection{Risk Evaluation}

%---------------------------------------

\subsection{Risk Treatment Options}

\subsubsection{Accept}

\subsubsection{Avoid}

\subsubsection{Transfer}

\subsubsection{Mitigate}

%---------------------------------------

\subsection{Healthcare Risk Register}

% Risk Register Table goes here

%%%%%%%%%%%%%%%%%%%%%%%%%%%%%%%%%%%%%%%%%%%%%%%%%%%%%

\section*{Chapter Summary}

Summarize the key findings from this chapter and briefly transition to the NIST Cybersecurity Framework discussed in the next chapter.